\chapter{A Concrete Naming Certificate for $\TwinSubstrateBridgeUp$}
\label{ch:concrete-instances-twinsubstratebridge-namecert}
\origin{human}

The $\TwinSubstrateBridgeUp$ packet realizes the twin-substrate architecture of \autoref{ch:twin-substrate-synthesis} as a finite bridge certificate. It sits beside the MetaCIC audit synthesis of \autoref{ch:concrete-instances-metacicauditsynthesis-namecert}, the GroundCompiler audit map of \autoref{ch:concrete-instances-groundcompilerauditmap-namecert}, and the audit-map frontier packet of \autoref{ch:concrete-instances-auditmapfrontierpacket-namecert}. Those sibling chapters provide the two substrate rows and the frontier boundary; this chapter records the bridge packet that lets a consumer read them together without identifying their audit ledgers.

\begin{definition}[TwinSubstrateBridge carrier]
\label{def:twin-substrate-bridge-carrier}
A $\TwinSubstrateBridgeUp$ carrier is a finite $\BHist$ packet
$$
B=(M,G,F,R,L,H,C,P,N)
$$
with the following visible rows:
\begin{enumerate}
\item $M$ is the MetaCIC audit-synthesis row consumed from \autoref{ch:concrete-instances-metacicauditsynthesis-namecert};
\item $G$ is the GroundCompiler audit-map row consumed from \autoref{ch:concrete-instances-groundcompilerauditmap-namecert};
\item $F$ is the audit-map frontier packet row consumed from \autoref{ch:concrete-instances-auditmapfrontierpacket-namecert};
\item $R$ is the bridge-route classifier separating MetaCIC-facing, GroundCompiler-facing, and frontier-facing reads;
\item $L$ is the two-sided bridge ledger recording which substrate row each read consumes;
\item $H$ stores componentwise $\hsame$ transport for $M,G,F,R,L$;
\item $C$ stores displayed $\Cont$ routes through the bridge rows;
\item $P$ records $\Pkg$ provenance over the sibling chapters, the route classifier, the ledger, and the structural rows;
\item $N$ is the local $\NameCert_{\TwinSubstrateBridgeUp}$ row.
\end{enumerate}
The carrier has no row identifying $M$ with $G$, no common audit equality, no compiler report promoted to MetaCIC proof evidence, and no MetaCIC derivation promoted to compiler acceptance.
\end{definition}
\leanchecked{BEDC.Derived.TwinSubstrateBridgeUp.TasteGate.TwinSubstrateBridgeTasteGate\_single\_carrier\_alignment}

\begin{theorem}[TwinSubstrateBridge NameCert obligations]
\label{thm:twin-substrate-bridge-namecert-obligations}
For every accepted $\TwinSubstrateBridgeUp$ carrier $B=(M,G,F,R,L,H,C,P,N)$, the local $\NameCert_{\TwinSubstrateBridgeUp}$ surface consists exactly of carrier admission, MetaCIC audit-synthesis readback through $M$, GroundCompiler audit-map readback through $G$, frontier readback through $F$, bridge-route classification through $R$, two-sided ledger exactness through $L$, componentwise $\hsame$ stability through $H$, displayed $\Cont$ replay through $C$, $\Pkg$ provenance through $P$, and local naming through $N$.
\end{theorem}
\leanchecked{BEDC.Derived.TwinSubstrateBridgeUp.TwinSubstrateBridgeCarrier\_namecert\_obligations}
\begin{proof}
Project the carrier in \autoref{def:twin-substrate-bridge-carrier}. The rows $M,G,F,R,L,H,C,P,N$ are the complete finite packet, so carrier admission and row readback are coordinate projections.

The route classifier $R$ is read before the ledger $L$. A MetaCIC-facing read is therefore sent to $M$ and may cite $F$ only as frontier context; a GroundCompiler-facing read is sent to $G$ and may cite $F$ only as frontier context. The ledger row records this choice rather than replacing either substrate row.

The structural rows preserve the same packet. Componentwise $\hsame$ transport in $H$ moves $M,G,F,R,L$ without merging them, $\Cont$ replay in $C$ exposes only displayed routes, and $P,N$ repack the same rows as provenance and local naming. Hence the NameCert surface is exactly the stated finite obligation list.
\end{proof}

\begin{theorem}[TwinSubstrateBridge carrier admission]
\label{thm:twin-substrate-bridge-carrier-admission}
Let $B=(M,G,F,R,L,H,C,P,N)$ be an accepted $\TwinSubstrateBridgeUp$ carrier. Carrier admission for $B$ is exactly the finite projection that exposes the MetaCIC audit-synthesis row $M$, the GroundCompiler audit-map row $G$, and the audit-map frontier row $F$ before any bridge route is consumed. No accepted carrier can expose $R,L,H,C,P,N$ as bridge evidence while omitting one of $M,G,F$.
\end{theorem}
\leanchecked{BEDC.Derived.TwinSubstrateBridgeUp.TwinSubstrateBridgeCarrier\_carrier\_admission}
\begin{proof}
Project the carrier of \autoref{def:twin-substrate-bridge-carrier}. The first three coordinates are the substrate rows that make the bridge meaningful: $M$ supplies the MetaCIC side, $G$ supplies the GroundCompiler side, and $F$ supplies the frontier boundary between them.

The later coordinates are dependent public rows. The route classifier $R$ can classify only requests between the displayed substrate sides, and the ledger $L$ can record only reads whose substrate side has already been exposed.

The structural rows $H,C,P,N$ preserve, replay, source, and name the same displayed packet. Since none of them manufactures a missing substrate coordinate, accepted carrier admission must expose $M,G,F$ before the bridge evidence can be read.
\end{proof}

\begin{theorem}[TwinSubstrateBridge replay stability]
\label{thm:twin-substrate-bridge-replay-stability}
For every accepted $\TwinSubstrateBridgeUp$ carrier $B=(M,G,F,R,L,H,C,P,N)$, componentwise transport through $H$ and displayed replay through $C$ preserve the route separation recorded by $R$ and $L$. A replayed MetaCIC-facing request still consumes $M$ before frontier context, a replayed GroundCompiler-facing request still consumes $G$ before frontier context, and a replayed bridge-facing request still compares the sides only through $F$ and $L$.
\end{theorem}
\leanchecked{BEDC.Derived.TwinSubstrateBridgeUp.TwinSubstrateBridgeCarrier\_replay\_stability}
\begin{proof}
Read the route classifier $R$ and ledger $L$ before replay. They determine whether a request is MetaCIC-facing, GroundCompiler-facing, or bridge-facing, and they record which substrate row the request may consume.

Transport through $H$ is componentwise. It moves the displayed rows $M,G,F,R,L$ along $\hsame$ without adding a row that identifies $M$ with $G$ or changes the route assigned by $R$.

Replay through $C$ follows the displayed continuation route after transport. Thus a MetaCIC-facing replay remains on the $M$ side, a GroundCompiler-facing replay remains on the $G$ side, and a bridge-facing replay reaches both sides only through the frontier and ledger rows. The separation is therefore stable under the structural rows of the carrier.
\end{proof}

\begin{theorem}[TwinSubstrateBridge ledger exactness]
\label{thm:twin-substrate-bridge-ledger-exactness}
Every public read of an accepted $\TwinSubstrateBridgeUp$ carrier $B=(M,G,F,R,L,H,C,P,N)$ factors through the route classifier $R$ and the two-sided ledger $L$ before a substrate row is consumed. Consequently, public consumers cannot read $M$, $G$, or $F$ as unclassified substrate data; the only exported substrate access is the route- and ledger-recorded access named by the local $\NameCert_{\TwinSubstrateBridgeUp}$ row.
\end{theorem}
\leanchecked{BEDC.Derived.TwinSubstrateBridgeUp.TwinSubstrateBridgeCarrier\_ledger\_exactness}
\begin{proof}
Start with any public read of the packet. The NameCert obligations in \autoref{thm:twin-substrate-bridge-namecert-obligations} expose the classifier and ledger as part of the finite public surface, so the read is first assigned a route by $R$ and then recorded by $L$.

Only after this routing step may the read consume a substrate row. The MetaCIC route reads $M$, the GroundCompiler route reads $G$, and the bridge route uses $F$ together with the ledger entry that records which side is being compared.

The remaining rows do not bypass this order. Componentwise transport in $H$, displayed replay in $C$, provenance in $P$, and local naming in $N$ preserve the same routed read. Hence every exported substrate access is exactly the access registered by $R,L$ and no unclassified read of $M,G,F$ is public.
\end{proof}

\begin{theorem}[TwinSubstrateBridge audit separation]
\label{thm:twin-substrate-bridge-audit-separation}
Every public consumer of an accepted $\TwinSubstrateBridgeUp$ carrier preserves the separation between the MetaCIC row $M$ and the GroundCompiler row $G$. A MetaCIC-facing consumer may read $M,F,R,L,H,C,P,N$; a GroundCompiler-facing consumer may read $G,F,R,L,H,C,P,N$; and a bridge-facing consumer may compare the two only through the displayed frontier row $F$ and ledger row $L$.
\end{theorem}
\begin{proof}
Begin with the route classifier $R$. It assigns each consumer request to the MetaCIC-facing, GroundCompiler-facing, or bridge-facing route before any audit row is exposed.

For the MetaCIC route, the only audit row supplied by the carrier is $M$. The frontier row $F$ and ledger row $L$ can record compiler-side context, but no row of the packet converts $G$ into proof evidence for $M$.

For the GroundCompiler route, the only compiler-facing audit row supplied by the carrier is $G$. The same frontier and ledger rows can record MetaCIC context, but no row converts $M$ into compiler acceptance evidence.

The bridge route reads both sides only through $F,L,H,C,P,N$. These rows display comparison context, transport, continuation, provenance, and naming; they do not add a common substrate equality. Thus every public consumer preserves the two audit ledgers as distinct coordinates of the same finite packet.
\end{proof}
\leanchecked{BEDC.Derived.TwinSubstrateBridgeUp.TwinSubstrateBridgeCarrier\_audit\_separation}

\begin{theorem}[TwinSubstrateBridge non-escape]
\label{thm:twin-substrate-bridge-nonescape}
An accepted $\TwinSubstrateBridgeUp$ carrier exports no hidden common substrate, host audit oracle, compiler proof checker, MetaCIC evaluator, collapsed equality between $M$ and $G$, or frontier-closing theorem. Every exported read factors through the finite rows $M,G,F,R,L,H,C,P,N$ and through the local $\NameCert_{\TwinSubstrateBridgeUp}$ surface.
\end{theorem}
\begin{proof}
Read the packet through \autoref{thm:twin-substrate-bridge-namecert-obligations}. The local NameCert row exposes only the displayed carrier, sibling readbacks, route classifier, ledger, transport, continuation, provenance, and local naming rows.

Apply \autoref{thm:twin-substrate-bridge-audit-separation}. Any request involving both substrates must pass through the frontier and ledger rows, so the packet cannot manufacture a common equality or a shared audit oracle.

The sibling chapters provide only their own named surfaces: MetaCIC synthesis, GroundCompiler audit mapping, and audit-map frontier accounting. Since the bridge carrier has no coordinate for a compiler proof checker, MetaCIC evaluator, hidden host oracle, or frontier-closing theorem, none of those objects can be exported by a consumer route.
\end{proof}
\leanchecked{BEDC.Derived.TwinSubstrateBridgeUp.TwinSubstrateBridgeCarrier\_nonescape}

\closureat{TwinSubstrateBridgeUp}{seedStr}

\begin{closurestatus}{\TwinSubstrateBridgeUp}
  \constructivestory{A $\TwinSubstrateBridgeUp$ packet is generated from a MetaCIC audit-synthesis row, a GroundCompiler audit-map row, an audit-map frontier row, a bridge-route classifier, a two-sided bridge ledger, componentwise $\hsame$ transports, displayed $\Cont$ routes, $\Pkg$ provenance, and a local $\NameCert$ row.}
  \theoryclosure{\seedClosure}
  \scopeclosed{\autoref{def:twin-substrate-bridge-carrier}, \autoref{thm:twin-substrate-bridge-namecert-obligations}, \autoref{thm:twin-substrate-bridge-audit-separation}, and \autoref{thm:twin-substrate-bridge-nonescape} bind the finite bridge carrier, local NameCert obligation surface, audit separation, and non-escape boundary.}
  \formalstatus{\unformalizedV}
  \bridgestatus{none}
  \origin{human}
  \notclaimed{This chapter does not close Lean verification, a public standard bridge, a common substrate equality, compiler proof checking, MetaCIC evaluation, frontier closure, or host audit automation.}
  \upgradepath{Formal closure requires checked carrier admission, substrate readback, route classification, ledger exactness, componentwise transport, continuation replay, provenance, naming, and non-escape for the same finite bridge packet.}
\end{closurestatus}
